% Left phi(s)=s^2 on[0,2], secant slope2 equals tangent slope at xi=1.
% Right h(z,t)=z^2+t-1=0, two branches, locally unique near(t,z)=(0,1).
\begin{tikzpicture}[font=\small,>={Latex[length=2mm]}]
 \begin{scope}[x=1.55cm,y=.7cm]
  \draw[->,black] (-.2,0)--(2.4,0) node[right] {$s$};
  \draw[->,black] (0,-1)--(0,4.8) node[left] {$\phi$};
  \draw[thick,domain=0:2.05,samples=40] plot(\x,{\x*\x});
  \draw[dashed,thick] (0,0)--(2,4);
  \draw[blue!70!black,thick,domain=.1:2.05] plot(\x,{2*\x-1});
  \fill (1,1) circle(1.4pt);\draw[dotted] (1,0)--(1,1);
  \node[below] at (1,0) {$\xi$};\node[above left] at (1,1) {parallel};
  \node at (1,5.2) {mean value theorem};
  \node[align=center] at (1,-2) {$\phi(s)=s^2$\\secant and tangent have equal slope};
 \end{scope}
 \begin{scope}[xshift=8cm,x=1.6cm,y=1.6cm]
  \draw[->,black] (-.6,0)--(1.65,0) node[right] {$t$};
  \draw[->,black] (0,-1.4)--(0,1.5) node[above] {$z$};
  \draw[blue!70!black,thick,domain=-.4:.999,samples=80] plot(\x,{sqrt(1-\x)});
  \draw[dashed,thick,domain=-.4:.999,samples=80] plot(\x,{-sqrt(1-\x)});
  \draw[dotted,thick] (0,1) circle(.25);
  \fill (0,1) circle(1.4pt);\node[anchor=west] at (.3,1.2) {local branch};
  \fill (1,0) circle(1.4pt) node[below right] {$(t,z)=(1,0)$};
  \node at (.4,2.25) {implicit function theorem};
  \node[align=center] at (.4,-1.95) {$h(z,t)=z^2+t-1=0$\\uniqueness is local near $(t,z)=(0,1)$};
 \end{scope}
\end{tikzpicture}
